
\section{Expander ensembles}
\label{sec:expander-enumerators}

Expand--Accumulate and Expand--Convolute use a sparse random matrix before a
recursive rate-one map.  Their published distance proofs bound the resulting
state process by a concentration inequality.  We instead keep its exact
input--output enumerator.  This section supplies the expander part of that
enumerator and proves the composition rule needed below.

Except in the final subsection, all vectors and matrices are binary.  We use the
convention that a binomial coefficient is zero outside its natural range.

\subsection{Independent Bernoulli entries}

Let $\dist^{\mathrm{Ber}}_{k,n,p}$ be the distribution on $k\times n$ matrices
whose entries are independent Bernoulli random variables with parameter
$p\in(0,1/2)$.  This is the ensemble used by the provable variants of
Expand--Accumulate and Expand--Convolute.

For $r\in[k]$, define
\begin{equation}
 q_r:=\frac{1-(1-2p)^r}{2}.
 \label{eq:bernoulli-expander-parity}
\end{equation}

\begin{lemma}[Bernoulli expander enumerator]
\label{lem:bernoulli-expander-enumerator}
For $r\in[k]$ and $u\in[0,n]$,
\begin{equation}
 \boxed{
 \enum^{\mathrm{Ber}}_{r,u}
   =
 \binom{k}{r}\binom{n}{u}q_r^u(1-q_r)^{n-u}.
 }
 \label{eq:bernoulli-expander-enumerator}
\end{equation}
\end{lemma}

\begin{proof}
Fix a message $\xv$ of weight $r$.  Each coordinate of $\xv\B$ is the parity
of $r$ independent Bernoulli variables.  The piling-up identity gives its
nonzero probability as $q_r$.  Different output coordinates use different
columns of $\B$ and are independent.  Thus $\xv\B$ has independent Bernoulli
coordinates with parameter $q_r$.  There are $\binom{k}{r}$ possible messages,
which proves \eqref{eq:bernoulli-expander-enumerator}.
\end{proof}

\subsection{Fixed column degree}

The original expression in this draft described a different ensemble.  We
record the normalized formula because fixed column degree may be useful for a
later implementation.  Let $\dist^{\mathrm{col}}_{k,n,d}$ sample each column
independently and uniformly from the binary vectors of weight $d$.  For a
message of weight $r$, put
\begin{align}
 O_{r,d}
   &:=\sum_{b\geq 0}\binom{r}{2b+1}\binom{k-r}{d-2b-1},
 &
 Z_{r,d}
   &:=\sum_{b\geq 0}\binom{r}{2b}\binom{k-r}{d-2b}.
 \label{eq:column-expander-odd-even}
\end{align}
Here $O_{r,d}$ and $Z_{r,d}$ count columns whose intersection with the message
support has odd and even cardinality, respectively.  In particular,
$O_{r,d}+Z_{r,d}=\binom{k}{d}$.

\begin{lemma}[Fixed-column expander enumerator]
\label{lem:column-expander-enumerator}
For $r\in[k]$ and $u\in[0,n]$,
\begin{equation}
 \boxed{
 \enum^{\mathrm{col}}_{r,u}
   =
 \binom{k}{r}\binom{n}{u}
 \left(\frac{O_{r,d}}{\binom{k}{d}}\right)^u
 \left(\frac{Z_{r,d}}{\binom{k}{d}}\right)^{n-u}.
 }
 \label{eq:column-expander-enumerator}
\end{equation}
\end{lemma}

\begin{proof}
Fix a message support of size $r$.  Each output coordinate is one precisely
when its sampled column has odd intersection with that support.  The columns
are independent, and their odd-intersection probability is
$O_{r,d}/\binom{k}{d}$.  Multiplication by $\binom{k}{r}$ counts all messages.
\end{proof}

The factors $\binom{k}{d}^{-n}$ in
\eqref{eq:column-expander-enumerator} are essential.  Omitting them counts
matrices instead of averaging over the matrix ensemble.

\subsection{Fixed row weight}

Concrete implementations often sample each row with an exact weight.  Let
$\dist^{\mathrm{row}}_{k,n,d}$ sample the $k$ rows independently and uniformly
from the binary vectors of length $n$ and weight $d$.  Define the binary
Krawtchouk polynomial
\begin{equation}
 K_a^{(n)}(j)
   :=
 \sum_{t\in\Zbb}(-1)^t\binom{j}{t}\binom{n-j}{a-t}.
 \label{eq:krawtchouk-definition}
\end{equation}

\begin{lemma}[Fixed-row expander enumerator]
\label{lem:row-expander-enumerator}
For $r\in[k]$ and $u\in[0,n]$,
\begin{equation}
 \boxed{
 \enum^{\mathrm{row}}_{r,u}
   =
 \binom{k}{r}\frac{\binom{n}{u}}{2^n\binom{n}{d}^{r}}
 \sum_{j=0}^{n}K_d^{(n)}(j)^r K_j^{(n)}(u).
 }
 \label{eq:row-expander-enumerator}
\end{equation}
\end{lemma}

\begin{proof}
Fix a message of weight $r$.  Its image is the sum of $r$ independent uniform
weight-$d$ rows.  Let $f_d$ be the indicator of the weight-$d$ Hamming slice.
The Fourier transform of $f_d$ at a character of weight $j$ is
$K_d^{(n)}(j)$.  Fourier inversion shows that the number of ordered row tuples
whose sum equals a fixed word of weight $u$ is
\[
 2^{-n}\sum_{j=0}^{n}K_d^{(n)}(j)^r K_j^{(n)}(u).
\]
Divide by $\binom{n}{d}^r$, multiply by the $\binom{n}{u}$ possible output
words, and then multiply by the $\binom{k}{r}$ possible messages.
\end{proof}

Formula \eqref{eq:row-expander-enumerator} is exact, but its alternating sum is
not a numerically stable evaluation method at large lengths.  A certified
implementation should use a recurrence, interval arithmetic, or an equivalent
coefficient representation.

The following recurrence gives such a representation using only nonnegative
terms.  For $u,v\in[0,n]$, set
\begin{equation}
 P_d(u,v)
 :=
 \frac{\binom{u}{j}\binom{n-u}{d-j}}{\binom{n}{d}},
 \qquad
 j:=\frac{u+d-v}{2},
 \label{eq:row-shell-transition}
\end{equation}
where the right-hand side is zero unless $j$ is an integer in
$[0,u]\cap[d-(n-u),d]$.  Define $\rho_{0,0}=1$ and
$\rho_{0,u}=0$ for $u>0$, and recursively define
\begin{equation}
 \rho_{r+1,v}:=\sum_{u=0}^{n}\rho_{r,u}P_d(u,v).
 \label{eq:row-shell-recurrence}
\end{equation}

\begin{lemma}[Positive fixed-row recurrence]
\label{lem:row-shell-recurrence}
For a fixed message of weight $r$, the probability that its image under
$\B\gets\dist^{\mathrm{row}}_{k,n,d}$ has weight $u$ is $\rho_{r,u}$.
Consequently,
\begin{equation}
 \enum^{\mathrm{row}}_{r,u}=\binom{k}{r}\rho_{r,u}.
 \label{eq:row-expander-positive-enumerator}
\end{equation}
\end{lemma}

\begin{proof}
Suppose that the sum of the first $r$ selected rows has support size $u$.
The next row is a uniform $d$-subset of $[n]$.  Its intersection size $j$
with the current support has probability
\[
 \frac{\binom{u}{j}\binom{n-u}{d-j}}{\binom{n}{d}}.
\]
The symmetric difference of the two supports has size $u+d-2j$.  This is
exactly the transition in \eqref{eq:row-shell-transition}.  Induction gives
\eqref{eq:row-shell-recurrence}, and multiplication by the number
$\binom{k}{r}$ of messages gives
\eqref{eq:row-expander-positive-enumerator}.
\end{proof}

Unlike \eqref{eq:row-expander-enumerator}, recurrence
\eqref{eq:row-shell-recurrence} has no cancellation.  Its state space after
$r$ rows is contained in $[0,\min\{rd,n\}]$, and every state has at most
$d+1$ outgoing transitions.

\subsection{Region-stratified left-regular rows}
\label{sec:region-left-regular-expander}

The fixed-row ensemble controls the total row weight but does not control
where the nonzero entries occur.  A stronger ensemble spreads every row over
the codeword.  Assume that $n=d\ell$, and partition the columns into $d$
consecutive regions $I_1,\ldots,I_d$, each of length $\ell$.  The distribution
$\dist^{\mathrm{reg}}_{k,n,d}$ samples independent indices
$J_{i,a}\gets[\ell]$ for $(i,a)\in[k]\times[d]$ and sets
\begin{equation}
 B_{i,(a-1)\ell+J_{i,a}}=1.
 \label{eq:region-left-regular-definition}
\end{equation}
All other entries are zero.  Every row therefore has weight $d$ and has
exactly one nonzero entry in each region.  We call this ensemble a
\emph{region-stratified left-regular expander}.  The shorter term
\emph{regular expander} refers to the same ensemble when the regional
constraint is clear.  The right degrees remain random.

For a region of length $\ell$, define $\rho^{(\ell)}_{0,0}=1$ and
$\rho^{(\ell)}_{0,u}=0$ for $u>0$.  For $r\geq0$, set
\begin{equation}
 \rho^{(\ell)}_{r+1,v}
 =\rho^{(\ell)}_{r,v-1}\frac{\ell-v+1}{\ell}
  +\rho^{(\ell)}_{r,v+1}\frac{v+1}{\ell},
 \label{eq:region-left-regular-shell-recurrence}
\end{equation}
where terms with indices outside $[0,\ell]$ are zero.

\begin{lemma}[Region-stratified shell law]
\label{lem:region-left-regular-shell-law}
Fix a message of weight $r$ and sample
$\B\gets\dist^{\mathrm{reg}}_{k,n,d}$.  The restrictions of $\xv\B$ to the
$d$ regions are independent.  In each region, its weight has distribution
$\rho^{(\ell)}_{r,\cdot}$, and conditional on weight $u$, its support is
uniform among the $u$-subsets of that region.  Consequently, the expected
number of weight-$r$ messages whose regional weights are
$(u_1,\ldots,u_d)$ is
\begin{equation}
 \boxed{
 \binom{k}{r}\prod_{a=1}^{d}\rho^{(\ell)}_{r,u_a}.
 }
 \label{eq:region-left-regular-enumerator}
\end{equation}
\end{lemma}

\begin{proof}
Within one region, the image is the parity of $r$ independent uniform unit
vectors.  If its current weight is $u$, the next unit vector decreases the
weight with probability $u/\ell$ and increases it with probability
$(\ell-u)/\ell$.  This gives
\eqref{eq:region-left-regular-shell-recurrence}.  Coordinate-permutation
symmetry gives the conditional uniformity.  The samples in distinct regions
are independent.  Multiplication by $\binom{k}{r}$ counts the messages.
\end{proof}

Two bounds make the regional enumerator practical outside the small-weight
range.  For $r\geq1$, define
\begin{equation}
 q^{\mathrm{reg}}_r
 :=\frac{1-\exp(-2r/\ell)}{2},
 \qquad
 \gamma_r:=\frac{e^r r!}{r^r}.
 \label{eq:regular-poisson-parameters}
\end{equation}

\begin{lemma}[Poissonized left-regular comparison]
\label{lem:regular-poisson-comparison}
Fix a weight-$r$ message $\xv$, sample
$\B\gets\dist^{\mathrm{reg}}_{k,n,d}$, and let
$\mathbf Z\in\Fbb_2^n$ have independent Bernoulli coordinates with parameter
$q^{\mathrm{reg}}_r$.  For every function
$F:\Fbb_2^n\to\mathbb R_{\geq0}$,
\begin{equation}
 \Ebb_{\B}[F(\xv\B)]
 \leq\gamma_r^d\Ebb_{\mathbf Z}[F(\mathbf Z)].
 \label{eq:regular-poisson-comparison}
\end{equation}
\end{lemma}

\begin{proof}
In each region, replace the $r$ unit-vector samples by independent Poisson
counts of mean $r/\ell$ at its $\ell$ coordinates.  Conditional on total count
$r$, these counts have the multinomial law of the original samples.  Without
conditioning, their parities are independent Bernoulli variables with
parameter $q^{\mathrm{reg}}_r$.  The regional total is Poisson with mean $r$,
so the conditioning event has probability
$e^{-r}r^r/r!=\gamma_r^{-1}$.  The regional totals are independent.  Dropping
their joint conditioning event from the numerator proves
\eqref{eq:regular-poisson-comparison}.
\end{proof}

Stirling's formula gives $\gamma_r=\Theta(\sqrt r)$.  The comparison penalty
is therefore polynomial in $r$ for fixed left degree $d$.

\begin{lemma}[Left-regular point-mass bound]
\label{lem:regular-point-mass}
For the same fixed message and every $\yv\in\Fbb_2^n$,
\begin{equation}
 \Pr_{\B}[\xv\B=\yv]
 \leq2^{d-n}\left(1+e^{-2r/\ell}\right)^n.
 \label{eq:regular-point-mass}
\end{equation}
\end{lemma}

\begin{proof}
In one region, the Fourier eigenvalue of the $r$-step unit-vector walk at a
character of weight $j$ is $(1-2j/\ell)^r$.  Fourier inversion bounds each
regional point probability by $2^{-\ell}\Lambda_{\ell,r}$, where
\[
 \Lambda_{\ell,r}:=
 \sum_{j=0}^{\ell}\binom{\ell}{j}
       \left|1-\frac{2j}{\ell}\right|^r.
\]
Pair $j$ with $\ell-j$.  For $j\leq\ell/2$, the inequality
$1-x\leq e^{-x}$ gives
\[
 \Lambda_{\ell,r}
 \leq2\sum_{j=0}^{\ell}\binom{\ell}{j}e^{-2rj/\ell}
 =2(1+e^{-2r/\ell})^\ell.
\]
The $d$ regions are independent.  Taking the $d$th power proves
\eqref{eq:regular-point-mass}.
\end{proof}

The regional constraint is the mathematical property needed below.  It can be
implemented by arranging the row samples as a bit matrix, transposing the
block and region indices, and applying independent permutations within the
regions.  We refer to that implementation as a block-to-region transpose;
the ensemble itself remains the region-stratified left-regular expander.

\subsection{Randomly labeled prime-field rows}
\label{sec:prime-field-expander}

We now give the prime-field analogue used later.  This is the weighted model,
not the addition-only model: every sampled expander edge receives an
independent uniform label in $\Fbb_p^*$.  This choice is the direct analogue
of the random nonzero weighters in nonbinary BAA and RAA codes.  It is also the
point at which the construction incurs one field multiplication per expander
edge.  If all edge labels are fixed to one, the distribution below depends on
the complete message-symbol composition, and none of the support-only formulas
in this subsection applies.

Put $s:=p-1$.  In the Bernoulli ensemble, an entry of $\B$ is zero with
probability $1-\theta$ and equals each element of $\Fbb_p^*$ with probability
$\theta/s$.  For a message support of size $r$, define
\begin{equation}
 q^{(p)}_r
 :=\frac{s}{p}\left(1-
       \left(1-\frac{p\theta}{s}\right)^r\right).
 \label{eq:prime-bernoulli-activation}
\end{equation}

\begin{lemma}[Labeled Bernoulli expander law]
\label{lem:prime-bernoulli-expander}
Fix $\xv\in\Fbb_p^k$ with support size $r$.  The coordinates of $\xv\B$
are independent.  Each coordinate is zero with probability $1-q^{(p)}_r$;
conditional on being nonzero, it is uniform on $\Fbb_p^*$.
\end{lemma}

\begin{proof}
For a nontrivial additive character $\chi$ of $\Fbb_p$, one active row
contribution $Z$ satisfies
\[
 \Ebb[\chi(Z)]
 =1-\theta+\frac{\theta}{s}\sum_{a\in\Fbb_p^*}\chi(a)
 =1-\frac{p\theta}{s}.
\]
The $r$ active rows are independent.  Fourier inversion on the additive group
of $\Fbb_p$ gives zero mass
$p^{-1}(1+s(1-p\theta/s)^r)$ and equal mass
$p^{-1}(1-(1-p\theta/s)^r)$ at every nonzero value.  Distinct columns use
disjoint samples and are therefore independent.
\end{proof}

The labeled regular ensemble retains the support sampling in
\eqref{eq:region-left-regular-definition} and independently labels every
chosen entry by a uniform element of $\Fbb_p^*$.  For one region, let
$\rho^{(p,\ell)}_{0,0}=1$.  Starting from terms outside $[0,\ell]$ equal to
zero, define
\begin{equation}
\begin{split}
 \rho^{(p,\ell)}_{r+1,v}
 ={}&\rho^{(p,\ell)}_{r,v-1}\frac{\ell-v+1}{\ell}
 +\rho^{(p,\ell)}_{r,v}\frac{v}{\ell}\frac{p-2}{p-1}\\
 &+\rho^{(p,\ell)}_{r,v+1}\frac{v+1}{\ell}\frac{1}{p-1}.
\end{split}
\label{eq:prime-regular-shell-recurrence}
\end{equation}

\begin{lemma}[Labeled regular shell law]
\label{lem:prime-regular-shell-law}
Fix a message whose support has size $r$.  The $d$ regional outputs are
independent.  In each region, the weight has distribution
$\rho^{(p,\ell)}_{r,\cdot}$; conditional on weight $u$, the output is uniform
among the $\binom{\ell}{u}s^u$ vectors of weight $u$ in $\Fbb_p^\ell$.
\end{lemma}

\begin{proof}
From a vector of weight $u$, a labeled unit step chooses a zero coordinate
with probability $(\ell-u)/\ell$ and raises the weight.  At a nonzero
coordinate, its uniform label cancels the current value with probability
$1/(p-1)$ and otherwise preserves the weight.  This proves
\eqref{eq:prime-regular-shell-recurrence}.  The walk is invariant under
coordinate permutations and independent nonzero coordinate scalings.  This
monomial symmetry is transitive on every $p$-ary Hamming shell, which proves
the conditional uniformity.
\end{proof}

For the labeled regular ensemble, put
\begin{equation}
 q^{(p),\mathrm{reg}}_r
 :=\frac{s}{p}\left(1-
       \exp\left(-\frac{pr}{s\ell}\right)\right).
 \label{eq:prime-regular-poisson-activation}
\end{equation}
The Poisson conditioning argument in
Lemma~\ref{lem:regular-poisson-comparison} is unchanged.  After removing the
conditioning, each coordinate is zero with probability
$1-q^{(p),\mathrm{reg}}_r$ and otherwise uniform on $\Fbb_p^*$.  Consequently,
for every nonnegative $F:\Fbb_p^n\to\mathbb R_{\geq0}$,
\begin{equation}
 \Ebb_{\B}[F(\xv\B)]
 \leq \gamma_r^d\Ebb_{\mathbf Z}[F(\mathbf Z)],
 \label{eq:prime-regular-poisson-comparison}
\end{equation}
where $\gamma_r=e^r r!/r^r$ and $\mathbf Z$ has the independent coordinate law
just displayed.

For later dense bounds, define
\begin{align}
 A_{p,\ell,r}
 &:=\left(1+s\exp\left(-\frac{pr}{s\ell}\right)\right)^\ell,
 \nonumber\\
 B_{p,\ell,r}
 &:=s^{\ell-r}
   \left(1+s^{-1}\exp\left(-\frac{pr}{\ell}\right)\right)^\ell.
 \label{eq:prime-regular-fourier-terms}
\end{align}

\begin{lemma}[Labeled regular point mass]
\label{lem:prime-regular-point-mass}
For every fixed support-$r$ message and every $\yv\in\Fbb_p^n$,
\begin{equation}
 \Pr_{\B}[\xv\B=\yv]
 \leq p^{-n}\bigl(A_{p,\ell,r}+B_{p,\ell,r}\bigr)^d.
 \label{eq:prime-regular-point-mass}
\end{equation}
\end{lemma}

\begin{proof}
A character of $\Fbb_p^\ell$ with support size $j$ is an eigenvector of one
labeled unit step with eigenvalue
$1-pj/(s\ell)$.  There are $\binom{\ell}{j}s^j$ such characters.  On the
nonnegative part of the spectrum,
$|1-pj/(s\ell)|^r\leq\exp(-prj/(s\ell))$, whose full binomial sum is
$A_{p,\ell,r}$.  On the negative part, substitute $h=\ell-j$ and use
\[
 \left|1-\frac{pj}{s\ell}\right|^r
 \leq s^{-r}\exp\left(-\frac{prh}{\ell}\right).
\]
The corresponding full binomial sum is $B_{p,\ell,r}$.  Fourier inversion
bounds a regional point by
$p^{-\ell}(A_{p,\ell,r}+B_{p,\ell,r})$.  Independence of the regions proves
\eqref{eq:prime-regular-point-mass}.
\end{proof}

\subsection{Composition without an explicit interleaver}

The expander output is exchangeable even though the constructions do not place
an interleaver after the expander.  This symmetry gives the same enumerator
composition rule as a uniform interleaver.

\begin{lemma}[Exchangeable serial composition]
\label{lem:exchangeable-serial-composition}
Let $\B\gets\dist_1$ be a random $k\times n$ matrix such that
$\B\pi$ and $\B$ have the same distribution for every $n\times n$ permutation
matrix $\pi$.  Independently sample an $n\times N$ matrix
$\C\gets\dist_2$.  Then
\begin{equation}
 \boxed{
 \enum^{\dist_1\dist_2}_{r,h}
   =
 \sum_{u=0}^{n}
 \frac{\enum^{\dist_1}_{r,u}\enum^{\dist_2}_{u,h}}
      {\binom{n}{u}}.
 }
 \label{eq:exchangeable-serial-composition}
\end{equation}
\end{lemma}

\begin{proof}
Fix a message $\xv$ and condition on $\hw{\xv\B}=u$.  Column-exchangeability
implies that $\xv\B$ is uniform on the weight-$u$ Hamming slice.  Averaging
over the independent matrix $\C$, the conditional probability of output
weight $h$ is therefore
$\enum^{\dist_2}_{u,h}/\binom{n}{u}$.  Sum over $u$, then sum over all messages
of weight $r$.
\end{proof}

All three expander ensembles above satisfy the hypothesis of
Lemma~\ref{lem:exchangeable-serial-composition}.

\subsection{Exact first moment for Expand--Accumulate}

Let $\A\in\Fbb_2^{n\times n}$ be the accumulator from
Section~\ref{sec:RA}.  Its input--output enumerator is
\begin{equation}
 \enum^{\acc}_{u,h}
   =
 \binom{h-1}{\lceil u/2\rceil-1}
 \binom{n-h}{\lfloor u/2\rfloor}
 \label{eq:accumulator-enumerator-recalled}
\end{equation}
for $u,h\geq 1$, together with
$\enum^{\acc}_{0,0}=1$ and $\enum^{\acc}_{u,0}=0$ for $u>0$.

For $u\in[0,n]$ and $L\in[0,n]$, define the accumulator slice tail
\begin{equation}
 Q_{n,L}(u):=
 \begin{cases}
  1,&u=0,\\[2mm]
  \displaystyle
  \sum_{j=\lceil u/2\rceil}^{\min\{u,L\}}
  \frac{\binom{L}{j}\binom{n-L}{u-j}}{\binom{n}{u}},&u>0.
 \end{cases}
 \label{eq:accumulator-slice-tail}
\end{equation}

\begin{lemma}[Positive accumulator slice tail]
\label{lem:accumulator-slice-tail}
If the accumulator input is uniform on the weight-$u$ slice, then its output
has weight at most $L$ with probability $Q_{n,L}(u)$.
\end{lemma}

\begin{proof}
For $u>0$, put $a=\lceil u/2\rceil$ and $b=\lfloor u/2\rfloor$.
Summing \eqref{eq:accumulator-enumerator-recalled} over $h\leq L$ gives
\[
 \sum_{h=1}^{L}\binom{h-1}{a-1}\binom{n-h}{b}.
\]
This sum counts weight-$u$ subsets by the location $h$ of their $a$-th
smallest element.  Counting the same subsets by the number $j$ of elements
in $[L]$ gives
\[
 \sum_{j=a}^{\min\{u,L\}}\binom{L}{j}\binom{n-L}{u-j}.
\]
Division by $\binom{n}{u}$ proves
\eqref{eq:accumulator-slice-tail}.  The case $u=0$ is immediate.
\end{proof}

\begin{theorem}[Exact Expand--Accumulate first moment]
\label{thm:expand-accumulate-first-moment}
Let $\B$ follow any expander ensemble in this section and set $\G:=\B\A$.
For $L\in[0,n]$, define
\begin{equation}
 \mu_{\mathrm{EA}}(L)
   :=
 \sum_{r=1}^{k}\sum_{u=0}^{n}\sum_{h=0}^{L}
 \frac{\enum^{\B}_{r,u}\enum^{\acc}_{u,h}}{\binom{n}{u}}.
 \label{eq:expand-accumulate-first-moment}
\end{equation}
Then $\mu_{\mathrm{EA}}(L)$ is the expected number of nonzero messages whose
encodings have weight at most $L$.  Consequently,
\begin{equation}
 \Pr\left[\exists\xv\neq 0:\ \hw{\xv\G}\leq L\right]
 \leq \mu_{\mathrm{EA}}(L).
 \label{eq:expand-accumulate-failure}
\end{equation}
In particular, the same expression upper-bounds the probability that the
generated code has minimum distance at most $L$.
\end{theorem}

\begin{proof}
Lemma~\ref{lem:exchangeable-serial-composition} and
\eqref{eq:accumulator-enumerator-recalled} give the expected number in
\eqref{eq:expand-accumulate-first-moment}.  The event in
\eqref{eq:expand-accumulate-failure} implies that this nonnegative count is at
least one.  Markov's inequality proves the probability bound.
\end{proof}

Unlike the earlier Markov-chain argument, Theorem~\ref{thm:expand-accumulate-first-moment}
does not replace the accumulator weight distribution by a concentration bound.
It applies equally to Bernoulli, fixed-column, and fixed-row expanders.

For the fixed-row ensemble, Lemmas~\ref{lem:row-shell-recurrence} and
\ref{lem:accumulator-slice-tail} turn
\eqref{eq:expand-accumulate-first-moment} into the positive formula
\begin{equation}
 \boxed{
 \mu_{\mathrm{EA}}^{\mathrm{row}}(L)
 =\sum_{r=1}^{k}\binom{k}{r}
   \sum_{u=0}^{n}\rho_{r,u}Q_{n,L}(u).
 }
 \label{eq:row-expand-accumulate-positive-first-moment}
\end{equation}
Thus the exact-row calculation requires only hypergeometric probabilities.
In particular, it avoids the alternating Krawtchouk sum in
\eqref{eq:row-expander-enumerator} in the message-weight range where direct
evaluation is most delicate.
